% !TeX program = pdflatex
% =============================================================================
% Vorlage: Arbeitsblatt (nicht bewertet, Übungsmaterial für den Unterricht)
% Kantonsschule KFR – Physik
%
% Verwendung:
%   1. Diese Datei nach arbeitsblaetter/<klasse>/<semester>/<thema>/ kopieren
%      und sinnvoll umbenennen, z. B.
%      arbeitsblaetter/5-klasse/1-hs/thermodynamik/waermeleitung-arbeitsblatt.tex
%   2. \Klasse/\Thema/\ArbeitsblattTitel anpassen (weiter unten, vor
%      \begin{document}). \Datum füllt sich automatisch mit dem aktuellen
%      Datum (\today); Lehrperson ist fest auf "Loris Keller" gesetzt.
%      Kein Tabellenraster mehr -- Klasse/Datum/Thema/Lehrperson stehen als
%      einfache Textzeile unter dem Titel-Banner.
%   3. Jedes Arbeitsblatt beginnt mit einem Lernziele-Block (\begin{infobox})
%      direkt nach der Kopfzeile, noch vor der ersten Aufgabe -- Stichpunkte
%      pro Arbeitsblatt anpassen, den Block selbst nicht entfernen.
%   4. Direkt danach folgt IMMER ein Theorie-Block (\begin{theoriebox}) --
%      die fachliche Grundlage (Definitionen, Formeln, Merksätze), die für
%      die nachfolgenden Aufgaben gebraucht wird, in eigener grüner Box.
%      Reihenfolge im Dokument ist fest: Lernziele (blau) -> Theorie (grün)
%      -> Aufgaben (blau, taskbox). Theorie-Inhalt vor den Aufgaben
%      ausformulieren, nicht als Lückenfüller stehen lassen.
%   5. Aufgaben mit \question[Punktzahl] ... ergänzen, Lösung jeweils in
%      \begin{solution}...\end{solution} setzen. \begin{taskbox}...\end{taskbox}
%      ist optional und rahmt eine Aufgabe farbig ein.
%   6. Lösungen ein-/ausblenden: \printanswers (Lehrerexemplar) bzw.
%      \noprintanswers (Schülerexemplar, Standard) weiter unten umschalten.
% =============================================================================
\documentclass[addpoints,11pt,a4paper]{exam}

\usepackage[T1]{fontenc}
\usepackage[utf8]{inputenc}
\usepackage[ngerman]{babel}
\usepackage{lmodern}
\usepackage[a4paper,margin=2.2cm,headheight=14pt]{geometry}
\usepackage{amsmath}
\usepackage{siunitx}
% Hausstil: zusammengesetzte Einheiten immer als echter Bruch (\frac), nicht
% als Inline-Schrägstrich oder \tfrac -- gilt global für jedes \si/\SI.
\sisetup{per-mode=fraction,fraction-command=\frac}
\usepackage{array,booktabs,tabularx,multirow}
\usepackage{enumitem}
\usepackage{xcolor}
\usepackage{colortbl}
\usepackage{tikz}
\usepackage{physics}
\usepackage{circuitikz}
\usepackage[most]{tcolorbox}
\usepackage{lastpage}

\usetikzlibrary{arrows.meta,calc,angles,quotes}

\usepackage{setspace}
\setlength{\parindent}{0pt}
\setlength{\parskip}{0.6em}
% Hausstil v2: etwas grosszügigerer Zeilenabstand als Standard-LaTeX (1.0),
% damit Fliesstext innerhalb der Boxen weniger gedrängt wirkt -- 1.08 ist
% bewusst dezent gewählt (nicht \onehalfspacing/1.5), um nicht unnötig
% Platz/Seiten zu kosten.
\setstretch{1.08}
% Etwas grosszügigerer Standardabstand zwischen Listenpunkten (gilt für
% jede \begin{itemize}, sofern nicht lokal überschrieben).
\setlist[itemize]{itemsep=0.4em}

% -----------------------------
% Farben (Corporate-Farbschema Physik KFR)
% -----------------------------
\definecolor{titleblue}{HTML}{183A6B}
\definecolor{accentblue}{HTML}{2E6F9E}
\definecolor{lightblue}{HTML}{EAF4FB}
\definecolor{verylightblue}{HTML}{F7FBFE}
\definecolor{lightgray}{HTML}{F4F6F8}
\definecolor{darkgray}{HTML}{4A4A4A}
% Theorie-Block: eigene Farbe (grün), damit er sich auf einen Blick von den
% blauen Lernziele-/Aufgaben-Boxen abhebt.
\definecolor{theorygreen}{HTML}{2E7D4F}
\definecolor{lighttheorygreen}{HTML}{EAF7EF}
% Gruppenarbeit-Box: eigene Farbe (Petrol/Teal), damit Partner-/Gruppenarbeit
% sich auf einen Blick von normalen Einzelaufgaben (blau, taskbox) und vom
% Theorie-Block (grün, theoriebox) abhebt. Siehe Regel weiter unten bei
% \newtcolorbox{groupbox}.
\definecolor{groupteal}{HTML}{1B7F72}
\definecolor{lightgroupteal}{HTML}{E8F5F3}

% -----------------------------
% Spaltentypen
% -----------------------------
\newcolumntype{Y}{>{\centering\arraybackslash}X}
\newcolumntype{P}[1]{>{\raggedright\arraybackslash}p{#1}}

% -----------------------------
% Boxen
% -----------------------------
% Hausstil v2: einheitlicher Abstand VOR/NACH jeder Box (Titelbox, Info-,
% Theorie-, Task-, Beispiel-, Gruppen-, Antwortbox), damit Boxen nicht
% direkt am umgebenden Fliesstext oder aneinander kleben. Als globale
% tcolorbox-Vorgabe gesetzt statt einzeln pro Box, damit sie für alle
% \newtcolorbox-Definitionen unten automatisch gilt.
\tcbset{before skip=10pt, after skip=10pt}
\newtcolorbox{titlebox}{
  enhanced,
  colback=titleblue,
  colframe=titleblue,
  boxrule=0pt,
  arc=3mm,
  left=3mm,right=3mm,top=3mm,bottom=3mm
}

% exam.cls rückt die questions-Liste um die Breite von "10." (plus \labelsep)
% ein (Platz für die Nummerierung "1.", "2." ...); wir messen dieselbe Distanz
% hier unabhängig nach.
\newlength{\aufgabenindent}
\settowidth{\aufgabenindent}{10.\hskip\labelsep}

% left/right-Padding bewusst grosszügiger als bei answerbox: exam.cls hängt
% die Nummerierung von \question bündig an den inneren Rand des
% Textbereichs -- bei zu knappem Padding sitzt diese Nummer direkt auf dem
% Rahmen der Box statt sichtbar innerhalb davon.
% infobox steht ausserhalb von questions (z.B. beim Lernziele-Block auf
% Dokumentebene) und braucht deshalb KEINE Korrektur -- \linewidth entspricht
% dort bereits \textwidth.
\newtcolorbox{infobox}{
  enhanced,
  breakable,
  colback=lightblue,
  colframe=accentblue,
  boxrule=0.7pt,
  arc=2mm,
  left=6mm,right=6mm,top=3.5mm,bottom=3.5mm
}

% theoriebox: wie infobox, aber grün und mit farbigem Titelbalken (analog zu
% taskbox), damit der Theorie-Block auf einen Blick von Lernziele (blau,
% ohne Titelbalken) und Aufgaben (blau, mit Titelbalken) unterscheidbar ist.
% Steht wie infobox ausserhalb von questions -- KEINE Einrückungskorrektur
% nötig (siehe Kommentar bei infobox oben).
% Regel: hervorgehobene/definierte Begriffe INNERHALB einer theoriebox
% immer mit \textbf{\emph{...}} setzen (fett UND kursiv), nicht mit blossem
% \emph{...} -- damit Fachbegriffe beim Durchlesen der Theorie sofort
% auffallen. Ausserhalb von theoriebox (beispielbox, Aufgabentext, ...)
% bleibt \emph{...} allein korrekt.
\newtcolorbox{theoriebox}[1][Theorie]{
  enhanced,
  breakable,
  colback=lighttheorygreen,
  colframe=theorygreen,
  title={#1},
  fonttitle=\bfseries,
  boxrule=0.7pt,
  arc=2mm,
  left=6mm,right=6mm,top=3.5mm,bottom=3.5mm
}

% taskbox steht direkt in der questions-Liste (vor dem ersten \question),
% also innerhalb von deren Einrückung. Empirisch (per Pixelvergleich mit der
% Kopfzeilen-Tabelle) verschiebt "enlarge left by" hier die GESAMTE Box
% (nicht nur ihren linken Rand) -- ein positiver Wert schiebt sie nach
% RECHTS, nicht nach links, entgegen der naheliegenden Lesart des
% Schlüsselnamens. Ein negativer Wert (-\aufgabenindent) schiebt die Box
% exakt um die Listen-Einrückung nach links; width=\textwidth stellt danach
% wieder die volle Breite her, sodass beide Kanten mit der Tabelle
% übereinstimmen (verifiziert per Pixelmessung, Abweichung <=1px @300dpi).
% Falls taskbox künftig auch ausserhalb von questions verwendet wird, muss
% diese Korrektur analog zu infobox entfernt werden.
\newtcolorbox{taskbox}[1][]{
  enhanced,
  breakable,
  colback=verylightblue,
  colframe=accentblue,
  title={#1},
  fonttitle=\bfseries,
  boxrule=0.7pt,
  arc=2mm,
  enlarge left by=-\aufgabenindent,
  width=\textwidth,
  left=6mm,right=6mm,top=3.5mm,bottom=3.5mm
}

% groupbox: Regel -- JEDE Aufgabe, bei der Schülerinnen/Schüler zu zweit
% oder in einer Gruppe zusammenarbeiten (Partnerarbeit, Gruppenarbeit,
% Diskussion zu zweit, ...), bekommt diese Box statt der blauen taskbox.
% Farbe (Petrol/Teal) ist bewusst weder blau (taskbox/infobox) noch grün
% (theoriebox), damit "hier braucht es eine Partnerin/einen Partner" auf
% einen Blick erkennbar ist. Strukturell identisch zu taskbox (gleiche
% Einrückungskorrektur, da auch groupbox innerhalb von questions steht).
\newtcolorbox{groupbox}[1][Gruppenarbeit]{
  enhanced,
  breakable,
  colback=lightgroupteal,
  colframe=groupteal,
  title={#1},
  fonttitle=\bfseries,
  boxrule=0.7pt,
  arc=2mm,
  enlarge left by=-\aufgabenindent,
  width=\textwidth,
  left=6mm,right=6mm,top=3.5mm,bottom=3.5mm
}

\newtcolorbox{answerbox}{
  breakable,
  colback=white,
  colframe=gray!55,
  boxrule=0.5pt,
  arc=1.5mm,
  left=2mm,right=2mm,top=2mm,bottom=2mm
}

% --- Lösungen ein-/ausblenden -----------------------------------------------
% Schülerexemplar (Standard): \noprintanswers
% Lehrer-/Lösungsexemplar:    \printanswers
\noprintanswers
% \printanswers

% --- Punkte bleiben intern, werden aber nicht gedruckt ----------------------
% Arbeitsblätter sind nicht benotet. \question[n]/\part[n] bekommen trotzdem
% eine Punktzahl n -- rein als interner Anhaltspunkt, der 1:1 an
% \Antwortraum{n} weitergereicht wird, um die Antwortraum-Grösse zu
% bestimmen. \pointformat{} (Aufruf, nicht \renewcommand!) unterdrückt die
% von der exam-Klasse sonst automatisch gedruckte Punktanzeige ("(n P.)")
% vollständig, ohne die interne Punktezählung (\numpoints) zu beeinflussen.
\pointformat{}

% --- Lösungsbox auf Deutsch, farblich abgesetzt -----------------------------
\renewcommand{\solutiontitle}{\noindent\color{titleblue}\textbf{Lösung:}\enspace}

% --- Kopf-/Fusszeile: Thema/Titel oben, Seitenzahl unten --------------------
\pagestyle{headandfoot}
\cfoot{\textcolor{darkgray}{Seite \thepage\ von \pageref{LastPage}}}

% --- Kopfzeile: Klasse / Datum / Thema / Titel ------------------------------
% Einheitliches Layout für Arbeitsblatt, Prüfung und Praktikum. Ohne Tabelle:
% Klasse/Datum/Thema/Lehrperson stehen als einfache Textzeile unter dem
% Titel-Banner.
\newcommand{\Kopfzeile}[4]{%
  \begin{titlebox}
    {\color{white}\sffamily\bfseries\Large #4}
  \end{titlebox}
  \noindent\textbf{Klasse:}~#1 \hfill \textbf{Datum:}~#2 \hfill \textbf{Thema:}~#3
    \hfill \textbf{Lehrperson:}~Loris Keller
  \lhead{\textcolor{darkgray}{#3}}
  \rhead{\textcolor{darkgray}{#4}}
}

% --- Antwortraum: ca. 2.5 cm Platz pro Punkt, als umrandetes Feld -----------
\newlength{\antwortraumlen}
\newcommand{\Antwortraum}[1]{%
  \pgfmathsetlength{\antwortraumlen}{#1*2.5cm}%
  \begin{answerbox}
    \rule{0pt}{\antwortraumlen}%
  \end{answerbox}%
}

% Werte für Kopfzeile zentral setzen
\newcommand{\Klasse}{5b}
\newcommand{\Datum}{\today}
\newcommand{\Thema}{Kinematik}
\newcommand{\ArbeitsblattTitel}{Arbeitsblatt: Kinematik}

\begin{document}

\Kopfzeile{\Klasse}{\Datum}{\Thema}{\ArbeitsblattTitel}

% --- Lernziele: fester Einstieg, vor der ersten Aufgabe ---------------------
\begin{infobox}
\textbf{Lernziele}\\
Am Ende dieser Einheit kann ich \dots
\begin{itemize}[leftmargin=1.2cm,itemsep=0.4em]
  \item \dots
  \item \dots
\end{itemize}
\end{infobox}

% --- Theorie: fachliche Grundlage für die Aufgaben, immer vor der ersten
%     Aufgabe. Definitionen/Formeln/Merksätze ausformulieren -- kein
%     Platzhalter, den man überspringen kann. ---------------------------------
\begin{theoriebox}
\dots
\end{theoriebox}

\begin{questions}

  \begin{taskbox}[Aufgabe 1]
  \question[4] Ein Zug fährt mit konstanter Geschwindigkeit
  $v = \SI{80}{\kilo\meter\per\hour}$.

  \begin{parts}
    \part[2] Rechnen Sie $v$ in $\si{\meter\per\second}$ um.
    \Antwortraum{2}

    \part[2] Wie weit fährt der Zug in $\SI{5}{\minute}$?
    \Antwortraum{2}
  \end{parts}

  \begin{solution}
    \textbf{a)} $v = \SI{80}{\kilo\meter\per\hour} = \frac{80\,000\,\si{\meter}}{3600\,\si{\second}} \approx \SI{22.2}{\meter\per\second}$

    \textbf{b)} $s = v \cdot t = \SI{22.2}{\meter\per\second} \cdot \SI{300}{\second} \approx \SI{6667}{\meter}$
  \end{solution}
  \end{taskbox}

\end{questions}

\end{document}
